%! Solving Rational Equations — Unit 4, Lesson 4.6 — Lesson Plan
\documentclass[10pt]{article}
\usepackage{algebra2-article}
\usepackage{algebra2-boxes}
\usepackage{graphicx}   % -article does NOT load graphicx; needed for thumbnails/screenshots
\graphicspath{{images/}}
\newcommand{\TallMath}[1]{$\displaystyle #1\rule[-1.4em]{0pt}{3.2em}$}
\newcommand{\CourseName}{Algebra 2}
\newcommand{\MeetingLength}{55 minutes}
\newcommand{\UnitNumberName}{Unit 4: Rational Functions \quad}
\newcommand{\LessonNumberName}{Lesson 4.6: Solving Rational Equations}

\begin{document}
\begin{center}
    {\Large\bfseries \CourseName \\
    {\small\bfseries \UnitNumberName \LessonNumberName}}
\end{center}

\begin{tcolorbox}[colback=forestbg, colframe=forest, boxrule=0.9pt, arc=2mm,
    left=2mm, right=2mm, top=1.5mm, bottom=1.5mm]
    \textbf{Primary Objective:} Students can solve an equation containing rational expressions with the
    same four steps every time: \textbf{factor} every denominator and write the \textbf{restriction
    list first}; \textbf{multiply every term} by the LCD --- both \emph{sides}, not just the fractions
    --- so the denominators cancel; \textbf{solve} the linear or quadratic equation that is left; and
    \textbf{check} each candidate against the restriction list, discarding any match as
    \textbf{extraneous} and verifying the survivors in the original. They can \emph{justify} why an
    extraneous solution appears (multiplying both sides by an expression that could be zero is not
    reversible), verify a solution set \emph{graphically} --- true solutions are $x$-intercepts,
    extraneous candidates are holes --- and \emph{model} a situation with a rational equation,
    interpreting the answer and distinguishing a solution rejected by the algebra from one rejected by
    the context.
    \\[2pt]
    \textbf{Standards (2023 VA SOL):} \textbf{A2.EI.4b} --- solve rational equations with real
    solutions containing factorable linear and quadratic expressions, algebraically and graphically.
    \textbf{A2.EI.4c} --- verify possible solutions algebraically, graphically, and with technology;
    explain the solution method and interpret solutions in context. \textbf{A2.EI.4d} --- justify why a
    possible solution might be extraneous. \textbf{A2.EI.4a} --- create an equation containing a
    rational expression to model a contextual situation. Applied: \textbf{A2.F.2h} and the 4.5 cancel
    test (the graphical check), \textbf{A2.EO.1a/b} (4.1--4.3 restrictions and the LCD). Builds on
    \textbf{A2.EO.3b} (3.2 factoring) and Unit 2's quadratic solving. Prerequisite for \textbf{5.4},
    where the same extraneous-solution logic returns for radical equations.
\end{tcolorbox}

\renewcommand{\arraystretch}{1.4}
\begin{skillbox}[Priority Ideas \& Skills]{goldbox}
\begin{tabularx}{\linewidth}{@{}X|X@{}}
\textbf{Priority Ideas \& Skills} & \textbf{Key Understandings (the ``why'')} \\[2pt]
\hline
\vspace{-6pt}
\begin{itemize}[leftmargin=1.1em, nosep, after=\vspace{2pt}]
  \item Tell an \emph{expression} from an \emph{equation}, and clear denominators only in the second.
  \item Factor every denominator and write the restriction list \emph{before} solving.
  \item Multiply every term by the LCD --- including terms with no denominator.
  \item Solve the resulting linear or quadratic equation with the toolkit already owned.
  \item Check every candidate against the restriction list; discard matches; verify survivors.
  \item Justify, in words, why an extraneous solution appears.
  \item Verify a solution set graphically: $x$-intercepts are solutions, holes are extraneous.
  \item Write a rational equation from a context, solve it, and interpret --- including rejections.
\end{itemize}
&
\vspace{-6pt}
\begin{itemize}[leftmargin=1.1em, nosep, after=\vspace{2pt}]
  \item \textbf{Clearing the denominators solves a \emph{different} equation. The check is what brings you back.}
  \item Multiplying both sides by a nonzero number is reversible; multiplying by an expression that could be zero is not. That single asymmetry is the whole phenomenon.
  \item An extraneous solution is never a random number --- it is always a value from the Step 1 list. That is what makes Step 4 fast.
  \item The converse fails: every extraneous solution is an excluded value, but most excluded values never turn up as candidates.
  \item The Hook proves the check is mandatory: two equations with \emph{identical} algebra and different answers.
  \item ``No solution'' is a complete answer --- and it has two distinct causes worth telling apart.
  \item Rejecting for \emph{context} is not the same as rejecting for \emph{extraneousness}; 4.7 leans on the difference.
\end{itemize}
\end{tabularx}
\end{skillbox}

\begin{skillbox}[{Vocabulary, Concepts, \& Theorems}]{sky}
\renewcommand{\arraystretch}{1.3}
\begin{tabularx}{\linewidth}{@{}l X@{}}
\textbf{Rational equation} & An equation with the variable in a denominator. Its solutions must lie in the domain of \emph{every} expression in it. \\
\textbf{Restriction list} & The excluded values, taken off the \emph{original} denominators (4.1). Written first, used last. \\
\textbf{Clearing the denominators} & Multiplying \emph{every term} by the LCD so all denominators cancel. Not the same as 4.3's combining. \\
\textbf{Candidate} & A number the cleared equation produced. It is not an answer until it survives the check. \\
\textbf{Extraneous solution} & A candidate that solves the cleared equation but not the original --- always one of the excluded values. \\
\textbf{Cross-multiplying} & Step 2 pre-cancelled, valid \emph{only} for a proportion (one fraction $=$ one fraction). \\
\textbf{Graphical check} & Move everything to one side and combine: \TallMath{\text{solutions}=x\text{-intercepts}; \quad \text{extraneous}=\text{holes}} \\
\end{tabularx}
\end{skillbox}

\begin{fixedskillbox}[Activate Prior Knowledge \& Spiral Review]{forestbg}
    The Warm-Up sets up two tools and one idea. \textbf{(1)} \emph{Restrictions} off the original
    denominator, four quick items --- this is the list students will check against all period, and the
    phrase to insist on is ``\emph{original}.'' \textbf{(2)} \emph{The LCD}, three items, including one
    whose terms include a plain number --- because today the LCD multiplies every term, and the term
    with no denominator is the one students drop. Say out loud that 4.3 used the LCD to \emph{build} a
    common denominator and today it is used to \emph{destroy} all of them: same object, opposite job.
    \textbf{(3)} \emph{A fact about equations}, and this is the one to debrief properly: start from the
    false statement $5=2$, multiply both sides by $3$ (still false, and reversible by dividing), then
    multiply both sides by $0$ and get $0=0$ --- true, and \emph{not} reversible. There is no algebra in
    it, so nobody can hide, and it is the entire justification for A2.EI.4d. Land one sentence before
    moving on: \emph{multiplying both sides by something that might be zero can create solutions that
    were never there.}
\end{fixedskillbox}

\begin{skillbox}[Hook]{forestbg}
    Put two equations on the board side by side (they are in the notes):
    $\frac{x^2}{x-2}=\frac{4}{x-2}$ and $\frac{x^2}{x-3}=\frac{4}{x-3}$. Have the class clear the
    denominator on each. Both collapse to \emph{the same equation}, $x^2=4$, with the same candidates
    $x=2$ and $x=-2$. \textbf{``So these two equations have the same solutions. Yes?''} Let them agree.
    Then fill the four-row check table together. Three rows come out true; one row --- equation (I) at
    $x=2$ --- comes out $\frac40=\frac40$. Ask directly: \textbf{``is that row true or false?''} The
    answer is neither, and that is the whole point of the lesson. A number that makes the equation
    \emph{undefined} was never a candidate for making it true. The Hook is chosen so the algebra is
    literally identical, which kills in advance the belief students will otherwise carry all period ---
    that a careful enough solver would not need to check. Post the conclusion and leave it up:
    \textbf{the algebra could not tell them apart; only the check could.}
\end{skillbox}

\begin{skillbox}[Lesson]{forestbg}
\begin{multicols}{2}
\textbf{1. Expression or equation?} Two minutes, and worth every one. Students just spent 4.3 building
common denominators, and under pressure they will ``clear'' the denominators of an \emph{expression}.
Put $\frac{6}{x}-\frac{2}{x-1}$ next to $\frac{6}{x}-\frac{2}{x-1}=1$ and name the difference: combining
keeps the denominator, clearing removes it, and only an equation has two sides to multiply. Post
\textbf{you may clear denominators only when there is an equals sign}, and expect to say it again during
the activity.

\textbf{2. The four steps.} Give them as a fixed order and hold to it: factor and restrict; multiply
every term by the LCD; solve; check. Step 1 is not bookkeeping --- it is the only thing that makes Step
4 possible, and it must be on the page before anyone touches the equation. Step 2 has one trap and it
is the biggest mechanical error of the day: the term with no denominator still gets multiplied. Write
$\cdot\,x(x-1)$ under \emph{every} term when you demonstrate.

\textbf{3. Why it happens (A2.EI.4d).} Cash in the Warm-Up. The LCD is not a number; it is an
expression, and at each excluded value it \emph{is} zero. So the multiplication is irreversible exactly
there, and the cleared equation can gain answers --- but only those. Say the consequence plainly and
write it: \textbf{an extraneous solution is never a random number; it is always on your Step 1 list.}
That is what makes Step 4 a comparison of two short lists rather than a second round of substitution.

\columnbreak

\textbf{4. Worked Example A --- $\frac{6}{x}-\frac{2}{x-1}=1$.} The case where nothing goes wrong. Run
all four steps out loud; candidates $2$ and $3$, restriction list $\{0,1\}$, no overlap, both verify.
Students need one clean run before they meet a failure, so that ``check'' does not read as ``expect
disaster.''

\textbf{5. Worked Example B --- $\frac{x}{x+1}+\frac{2}{x-1}=\frac{2}{x^2-1}$.} Factor the last
denominator first. Candidates $0$ and $-1$; $-1$ is on the list, so it goes --- \emph{without}
substituting anything. Then the graphical check, which is where 4.5 pays its debt: combine everything
onto one side and the result is $\frac{x(x+1)}{(x-1)(x+1)}=\frac{x}{x-1}$, a cancelling factor and a
surviving one. Yesterday's cancel test sorts today's candidates. Say the pairing and box it:
\textbf{true solutions are $x$-intercepts; extraneous candidates are holes.} Anyone with graphing
technology can now verify a solution set in ten seconds (A2.EI.4c).

\textbf{6. Two last things.} Cross-multiplying is Step 2 pre-cancelled and applies \emph{only} to a
proportion --- present it that way or students will reach for it on three-term equations and lose a
term. And ``no solution'' is a complete answer: run
$\frac{1}{x-2}+\frac{1}{x+2}=\frac{4}{x^2-4}$, whose only candidate is the excluded $2$. Note for
yourself that there is a second route to ``no solution'' --- the cleared equation being itself never
true --- and that homework 2(d) versus 2(f) puts both on one page.
\end{multicols}
\end{skillbox}

\begin{skillbox}[Explicit Instruction: Solve, Then Check]{forestbg}
\begin{tabularx}{\linewidth}{@{}p{0.44\linewidth} X@{}}
\vspace{0pt}
\textbf{Post these. Always in this order.}
\begin{enumerate}[leftmargin=1.3em, nosep, after=\vspace{2pt}]
  \item \textbf{Factor \& restrict.} Every denominator factored; the excluded values written at the top of the page.
  \item \textbf{Multiply by the LCD.} \emph{Every term}, including any term with no denominator.
  \item \textbf{Solve} the linear or quadratic equation that is left.
  \item \textbf{Check.} Candidate on the list $\Rightarrow$ extraneous, discard. Verify the survivors in the \emph{original}.
\end{enumerate}

\vspace{3pt}
\textbf{The sentence:} \emph{clearing the denominators solves a different equation --- the check is what
brings you back.}
&
\vspace{0pt}
\textbf{Nothing thrown away.} \; $\dfrac{6}{x}-\dfrac{2}{x-1}=1$, \; $x\neq0,1$
\[ 6(x-1)-2x=x(x-1) \;\Rightarrow\; x^2-5x+6=0 \;\Rightarrow\; x=2,\,3 \]
\centerline{\small Neither is on $\{0,1\}$: \; solution set $\{2,3\}$.}

\vspace{4pt}
\textbf{One thrown away.} \; $\dfrac{x}{x+1}+\dfrac{2}{x-1}=\dfrac{2}{x^2-1}$, \; $x\neq1,-1$
\[ x(x-1)+2(x+1)=2 \;\Rightarrow\; x(x+1)=0 \;\Rightarrow\; x=0,\,-1 \]
\centerline{\small $-1$ is on the list: \; extraneous. \; Solution set $\{0\}$.}

\vspace{4pt}
\textbf{The two questions students must answer aloud:} \emph{did every term get multiplied?} and
\emph{is this candidate on my list?}
\end{tabularx}
\end{skillbox}

\begin{skillbox}[Active Monitoring]{redbox}
Circulate during the guided practice and Tier work. Watch for and correct:
\begin{itemize}[leftmargin=1.2em, nosep]
  \item \textbf{no restriction list on the page} --- the single most expensive omission today, because it makes Step 4 impossible rather than merely slow;
  \item \textbf{restrictions read off the \emph{simplified} form} --- 4.1's signature error, which today stops being a mislabeled feature and becomes a wrong answer;
  \item \textbf{only the fractions multiplied by the LCD}, leaving the constant term untouched (homework 5 is built on exactly this);
  \item \textbf{a lost distributed minus sign} when a subtracted fraction is cleared --- 4.3's lost parenthesis, still in force;
  \item \textbf{cross-multiplying a three-term equation} --- a term simply vanishes;
  \item \textbf{clearing the denominators of an \emph{expression}} that was only meant to be combined;
  \item \textbf{a negative or fractional answer rejected on reflex} --- being ugly is not being extraneous (activity Tier R, 2(d));
  \item \textbf{``extraneous'' used as a synonym for ``rejected''} --- a context rejection is a different thing (activity Tier E, homework 6);
  \item \textbf{``no solution'' refused}, and the whole problem restarted in search of a number.
\end{itemize}
Cold-call prompts: ``What is your restriction list? Show me where it is written.'' ``Did \emph{every}
term get multiplied --- including that one?'' ``Is that candidate on your list? Then what?'' ``Why did
that extra answer appear at all?'' ``Where would this solution be on a graph? Where would the
extraneous one be?'' ``You rejected that one --- was it the algebra or the situation that made you?''
\end{skillbox}

\begin{skillbox}[Group Work \& Differentiation]{redbox}
\textbf{Solve, Then Check} activity (tiered; mirrors the notes).
\begin{multicols}{3}
\textbf{Tier R --- Remediate.} One step at a time. A four-row \textbf{restrictions-and-LCD} table (no
solving); four \textbf{solution-or-not} items with the candidate already given; then \textbf{twins} ---
$\frac{x}{x-4}=\frac{4}{x-4}+3$ and $\frac{x}{x-4}=\frac{2}{x-4}+3$, which differ by one digit and come
out ``no solution'' and $x=5$. That pair is the tier's whole argument: nothing in the \emph{process}
predicted which. Watch item 2(d), where the valid answer is $x=-1$: students reject negatives on
reflex, and here the numbers to reject are positive.

\columnbreak
\textbf{Tier A --- Approaching.} \emph{Part 1:} a full solve of
$\frac{x}{x-2}+\frac{2}{x+2}=\frac{8}{x^2-4}$ --- candidates $-6$ and $2$, the latter extraneous.
\emph{Part 2:} the \textbf{graphical check} (A2.EI.4b/c) on $\frac{x^2}{x+3}=\frac{9}{x+3}$, whose
difference is the \emph{line} $y=x-3$ with the point at $x=-3$ removed: the solution is the
$x$-intercept $(3,0)$ and the extraneous candidate is the hole $(-3,-6)$, pre-drawn. \emph{Part 3:} an
\textbf{error analysis with no error in it} --- every line of the student's algebra on
$\frac{2}{x-5}+\frac{3}{x}=\frac{10}{x^2-5x}$ is correct and the answer is still wrong, because Step 4
never happened. Expect a group to insist there must be an arithmetic slip.

\columnbreak
\textbf{Tier E --- Extension.} \emph{Part 1:} \textbf{build the model} (A2.EI.4a) --- Dev takes $3$
hours longer than Mia, together $2$ hours, giving $\frac1x+\frac{1}{x+3}=\frac12$ and candidates $3$ and
$-2$. \emph{Neither is extraneous}; $-2$ is rejected because a job cannot take negative hours. Do not
let ``extraneous'' become a synonym for ``rejected.'' \emph{Part 2:} \textbf{make one on purpose} ---
design an equation whose only candidate is extraneous, then one with a valid answer and an extraneous
one. \emph{Part 3:} \textbf{can clearing ever \emph{lose} a solution?} On
$\frac{x^2}{x-1}=\frac{x}{x-1}$, a student who divides $x^2=x$ by $x$ discards the only real solution
($x=0$) and keeps the extraneous one ($x=1$) --- both failure modes in one line.
\end{multicols}
\end{skillbox}

\begin{skillbox}[Individual Work \& Assessment]{redbox}
\textbf{Exit Ticket} (independent, no notes): \textbf{(1)} a full solve of
$\frac{x}{x-1}+\frac{1}{x}=\frac{1}{x^2-x}$ --- factor $x^2-x=x(x-1)$, restrictions $\{0,1\}$, cleared
equation $x^2+(x-1)=1$, so $x^2+x-2=0$ and $(x+2)(x-1)=0$; $x=1$ is extraneous and the solution set is
$\{-2\}$. Two things to watch: a student who factors correctly but lists only $x\neq1$ has lost half the
list and still gets this item right, and a student who rejects $-2$ for being negative has imported a
rule from the modeling problems. \textbf{(2)} the \textbf{A2.EI.4d justification} in writing, required
to use the words \emph{denominator} and \emph{extraneous}; grade the \emph{reason}, not the verdict ---
``it doesn't check'' restates the answer and earns nothing, while naming the zero denominator in the
original \emph{and} connecting the extra root to the LCD step earns full credit. \textbf{(3)} an
\textbf{SOL-style multiple-choice} item on $\frac{x}{x+3}+\frac{3}{x-3}=\frac{18}{x^2-9}$, which clears
to $x^2=9$ so that \emph{both} candidates are excluded; the answer is \textbf{no solution}, and each
distractor keeps one candidate or both --- a student choosing ``$\{3,-3\}$'' did the algebra perfectly
and skipped Step 4. Sort into ``got it / no list / list unused / rejected for the wrong reason.'' The
middle two share one fast fix: require the list at the top of the page. \textbf{Watch the last pile
before 4.7}, where rejecting a solution \emph{for context} becomes the correct move and the two reasons
must stay separate.
\end{skillbox}

\begin{skillbox}[Reinforcement \& Extension]{goldbox}
\textbf{Homework:} four \textbf{restriction-and-LCD} items (no solving); six \textbf{solves} covering
every outcome --- two valid answers, one valid answer, and both flavors of ``no solution,'' with
\textbf{2(d) versus 2(f)} explicitly compared (in (d) a real candidate was thrown out as extraneous; in
(f) the cleared equation collapses to $0=4$ and there was never a candidate at all); three
\textbf{check-only} items; a \textbf{graph to read} that belongs to problem 2(e) --- the combined
function is $\frac{x+3}{x+1}$ with an $x$-intercept at $(-3,0)$, a hole at $(1,2)$, and a wall at
$x=-1$, so part (d) can ask the question the lesson has been circling: $-1$ is excluded and was
\emph{never} a candidate, which means \textbf{every extraneous solution is an excluded value, but not
every excluded value is extraneous}; an \textbf{error analysis} on the other big error of the day
(multiplying the fractions instead of both sides of $\frac{10}{x}-\frac{4}{x-1}=1$, giving
$x=\frac{11}{6}$ instead of $x=2,5$); and a \textbf{river-current model} (A2.EI.4a/c) --- $6$ miles up
and $6$ back at $4$ mph paddling, $4$ hours total, giving $\frac{6}{4-c}+\frac{6}{4+c}=4$ and
$c=\pm2$, where \emph{neither} candidate is extraneous and $c=-2$ is rejected by context, and where the
restriction $c\neq4$ is a 4.5 vertical asymptote that means the paddler exactly matches the river.
\textbf{Extension:} an equation whose restriction list contains $2$ though $2$ is never a candidate,
and the one-parameter family $\frac{x}{x-2}=\frac{k}{x-2}$, which has no solution for exactly one value
of $k$ --- the Hook, generalized. \textbf{Preview:} Lesson 4.7, \emph{Direct, Inverse \& Joint
Variation} --- the unit's modeling capstone, where $y=\frac kx$ goes to work and rejecting an answer
because the \emph{situation} forbids it becomes routine.
\end{skillbox}
\end{document}
